\section{Prise en compte du RGPD et limites avant mise en production}

Lynx Immo manipule des données relatives aux comptes utilisateurs, à leur progression, à leur profil et à leurs abonnements. Le \terme{rgpd} ne peut donc pas être réduit à une exigence documentaire ajoutée en fin de projet. Les choix d'architecture et d'accès aux données doivent contribuer à limiter l'exposition des informations et à distinguer les usages selon les rôles. Le texte réglementaire constitue ici un cadre de référence \parencite{rgpd-text}, mais le projet étudiant n'a pas fait l'objet d'une validation juridique complète ni d'un audit de conformité.

Plusieurs mesures concrètes ont néanmoins été intégrées dans la réalisation. L'authentification est centralisée avec Supabase Auth et l'accès aux tables exposées est contrôlé par des politiques \terme{rls}. Le principe retenu est qu'un client ne doit pas obtenir un accès général à la base puis filtrer lui-même les données : les règles d'accès sont appliquées au niveau de PostgreSQL selon l'identité et le rôle. Les fonctions serveur sont utilisées lorsque des traitements ne doivent pas reposer uniquement sur le client. Les exigences non fonctionnelles du projet identifient par ailleurs la sécurité, la confidentialité, la gestion des accès et la protection des données comme des dimensions à suivre ; leur état détaillé est reproduit en annexe~\ref{ann:ipf-exigences-non-fonctionnelles}.

La séparation entre l'application apprenant et l'administration répond également à des besoins différents. Les écrans administrateurs donnent accès à des fonctions de gestion qui ne sont pas nécessaires à un utilisateur standard ; leur accès doit donc être associé à des rôles spécifiques. De la même manière, RevenueCat est utilisé pour le cycle de vie des abonnements mobiles afin d'éviter de reconstruire localement toute la logique liée aux achats des plateformes. Cette délégation ne supprime pas la responsabilité sur les données échangées avec le service tiers : elle impose au contraire de documenter précisément les informations transmises et leur finalité.

Avant une mise en production commerciale, plusieurs éléments resteraient à formaliser avec le responsable de traitement et, si nécessaire, le conseil juridique ou le délégué à la protection des données : registre des traitements, finalités et bases légales, durées de conservation, procédure d'exercice des droits, suppression ou anonymisation des comptes, politique de confidentialité, gestion des sous-traitants, analyse des transferts éventuels hors Union européenne et procédure de traitement d'un incident. Une revue des données réellement collectées serait également nécessaire afin de supprimer celles qui ne sont pas utiles au service rendu.

Cette distinction est importante pour C1.6 et C2.4 : le projet montre une prise en compte technique de la confidentialité et du contrôle d'accès, mais je ne présente pas l'application comme juridiquement conforme sur la seule base de ces mécanismes. La conformité résulterait d'un ensemble plus large associant architecture, procédures, documentation, gouvernance des données et validation organisationnelle.